\section*{Chapter \ref{ch:mx:transaction-cost-analysis} --- Transaction-Cost Analysis}

\begin{solution}{exo:mx:transaction-cost-analysis:1}
Delay $800\times0.02=16$, execution $800\times0.03=24$, opportunity $200\times0.10=20$, fees $800\times0.003=2.4$ dollars: 62.4 dollars, or 6.24 cents
per share of the 1\,000-share target.
\end{solution}

\begin{solution}{exo:mx:transaction-cost-analysis:2}
The VWAP report says it beat the benchmark by a cent a share ($100.05-100.06$); the arrival report says it cost 3 cents a filled share, and with delay,
opportunity and fees 6.24 cents a share of the order.
\end{solution}

\begin{solution}{exo:mx:transaction-cost-analysis:3}
$-0.5+1.04\times10.0\times\sqrt{0.1}=2.79$ ticks a share.
\end{solution}

\begin{solution}{exo:mx:transaction-cost-analysis:4}
The desk sent the urgent algorithm the larger orders (7\,039 shares on average against 4\,126) and more of the volatile days' orders; both raise the
cost whatever the algorithm, so part of the raw 3.36 ticks is the orders' difficulty, not the algorithm.
\end{solution}

\begin{solution}{exo:mx:transaction-cost-analysis:5}
Reversion is measured after the last fill. The patient algorithm's last fills are often \omterm{def:mx:child-order-placement:risk}{clean-up trades} that cross the spread at the end of a slice,
after five minutes of resting pressure; the price then gives back that push (3.51 ticks after two minutes). The urgent algorithm's last slices are the
smallest of its front-loaded schedule, and its large early trades have already reverted during the window.
\end{solution}

\begin{solution}{exo:mx:transaction-cost-analysis:6}
When $1.89<\lambda(5.98^2-4.60^2)=14.6\lambda$, that is $\lambda>0.13$ per tick.
\end{solution}

\begin{solution}{exo:mx:transaction-cost-analysis:7}
2.32 ticks, with a 95\% interval from 0.79 to 3.85. Participation captures most of the size effect, but not the volatile days that the urgent
algorithm also received more often, so part of the bias remains (the truth is 1.89); the interval is narrower because the estimate uses fewer
controls.
\end{solution}

\begin{solution}{exo:mx:transaction-cost-analysis:8}
The interval VWAP includes the algorithm's own trades and the price they moved: it drifts with the order. The same orders cost 2.42 ticks a filled share
against the arrival price.
\end{solution}

\begin{solution}{pb:mx:transaction-cost-analysis:1}
\textbf{1.} The arrival price, the VWAP over the order's interval, the day's VWAP, the close.
\textbf{2.} Patient 2.42, $-0.29$, 1.90, 2.51; urgent 3.66, 1.40, 3.39, 4.10 ticks a filled share.
\textbf{3.} Its window's VWAP contains its own trades and their impact.
\textbf{4.} The arrival (or decision) price: fixed before the order touches the market, so it measures the whole cost of trading.
\textbf{5.} Half-spread plus $\eta\sigma\sqrt{Q/V}$; $\eta=1.04$ (standard error 0.18) with $-0.5$ tick of spread.
\textbf{6.} 2.47 predicted against 2.73 realised on average; a correlation of 0.25 order by order.
\textbf{7.} Delay, spread, impact, timing, opportunity and fees add up to Perold's shortfall because impact and timing split the execution cost exactly.
\textbf{8.} Patient: delay $-0.35$, spread $-0.65$, impact 2.41, timing 0.59, opportunity 0.10, fees $-0.20$, total 1.90; urgent: $-0.15$, 2.43, 1.13,
$-0.07$, 0.16, 0.29, total 3.79.
\textbf{9.} The price gives back 3.51 ticks (patient) and 1.06 (urgent) within two minutes: temporary impact, no leakage.
\textbf{10.} 53.5\% of orders to the urgent algorithm, more of them large or on volatile days (7\,039 shares on average against 4\,126).
\textbf{11.} 3.36 ticks, standard error 0.76.
\textbf{12.} Cost on the algorithm, the regime's volatility, participation, the square-root term and its interaction with the algorithm; clustered
by day because the day's market is common to its orders.
\textbf{13.} \emph{Named result}: urgent minus patient, adjusted for difficulty, is 2.00 ticks a share with a 95\% interval from 0.32 to 3.68, against
a raw difference of 3.36 (standard error 0.76); the truth, from working every order with both, is 1.89 (0.25).
\textbf{14.} Close: the adjustment removed most of the bias and the interval covers it. A real order is worked once; its other outcome is never seen.
\textbf{15.} Patient 0.77 ticks under the prediction (0.56), urgent 1.50 over it (0.46).
\textbf{16.} Everything to the patient algorithm: 2.17 ticks a share on the last 50 days against 3.24 for the desk's routing.
\textbf{17.} Its cost varies less (standard deviation 4.60 against 5.98 ticks): a manager penalising variance enough prefers it.
\textbf{18.} With orders assigned at random, the raw difference is unbiased and adjustment only narrows the interval.
\textbf{19.} Hard orders explain about 1.4 of the 3.4 ticks; the rest, about 2 ticks, is the algorithm, measured within an interval that a year of
such orders would narrow.
\textbf{20.} The same benchmark, fixed before trading, and an adjustment for how hard each order was.
\end{solution}

\begin{solution}{iq:mx:transaction-cost-analysis:1}
VWAP measures how the order's prices compared with the market's during the order, including the order's own push; arrival measures what the fund paid
against the price when it decided. The 20 basis points are impact and drift; the VWAP beat says the algorithm tracked the day.
\iqlookfor{Which benchmark moves with the order.}
\end{solution}

\begin{solution}{iq:mx:transaction-cost-analysis:2}
Spread against the mid just before each fill; impact from a pre-trade model (or from the market's move beta-adjusted to an index or sector), timing as
the remainder; or match similar orders on the opposite side to cancel the market's move.
\iqlookfor{A model or a control for the market's move.}
\end{solution}

\begin{solution}{iq:mx:transaction-cost-analysis:3}
The number of orders and the noise of each (standard errors clustered by day or stock), the difficulty of each broker's orders (size, volatility,
\omterm{def:mx:the-almgren-chriss-framework:urgency}{urgency}) and how orders were assigned; ideally randomised routing.
\iqlookfor{Standard errors, difficulty, assignment.}
\end{solution}

\begin{solution}{iq:mx:transaction-cost-analysis:4}
A price that gives back the order's push after completion shows temporary impact; one that keeps going shows information. Compare post-completion
drift with a pre-trade expectation, by broker and venue; drift that starts before the order's large child trades suggests leakage.
\iqlookfor{Sign conventions; timing relative to the order.}
\end{solution}

\begin{solution}{iq:mx:transaction-cost-analysis:5}
Parent: decision time and price, arrival, instructions, target, limit, algorithm and parameters, cancellations. Child: send, acknowledgement and fill
times, venue, price, quantity, fees, liquidity flag; market data around each (mid, spread, depth, volume) at microsecond time stamps.
\iqlookfor{Timestamps; the decision price; child-level venue data.}
\end{solution}

\begin{solution}{iq:mx:transaction-cost-analysis:6}
Ask for the universe's composition and the vendor's difficulty model, check whether the client's orders were routed differently (larger, more urgent,
more volatile names), rerun the comparison on the client's own orders adjusted with its own pre-trade model, and look at the interval.
\iqlookfor{Composition; adjustment; uncertainty.}
\end{solution}
